\chapter{Closing Remarks and Outlook}
\label{ch:conclusions}


The work presented in this dissertation is unified by a single underlying aspiration: to use the Universe itself as a precision laboratory for fundamental physics. Cosmological observables --- the anisotropies of the CMB, the large-scale distribution of matter, and the 21-cm signal from cosmic dawn --- each carry the imprints of physical processes that are inaccessible to any terrestrial experiment, from the dynamics of the inflationary epoch to the microphysics of dark matter and the free-streaming nature of neutrinos. Each of the three Parts of this dissertation has developed and deployed a distinct set of theoretical tools and observational strategies in service of this program, addressing different regimes of cosmic history and different classes of BSM physics. Taken together, they reflect a common methodology of identifying pristine observational signatures, developing the theoretical and computational infrastructure needed to extract them, and using the resulting constraints to probe --- and in some cases, rule out --- regions of BSM parameter space that would otherwise be out of reach.

\medskip

Part~\ref{part:warm-inflation} established warm inflation as a theoretically coherent and observationally distinctive framework, characterized by its own viability conditions, computational tools, and novel phenomenological signatures with no analogue in standard (cold) inflation. Part~\ref{part:dark-universe} probed the nature of dark matter and dark energy through a combination of cosmological and terrestrial observations, with the common lesson that the dark sector is accessible from multiple complementary angles and that the interpretation of even marginal observational evidence is inseparable from the theoretical assumptions brought to the data. Part~\ref{part:neutrinos} elevated a single observable --- the neutrino-induced phase shift --- into a pristine probe of neutrino interactions and free-streaming radiation more generally, across multiple cosmological datasets from the CMB to the 21-cm signal from cosmic dawn.


\medskip

The coming decade promises to be transformative for precision cosmology and its implications for fundamental physics. In the near term, SO will deliver CMB measurements of unprecedented sensitivity, tightening constraints on the inflationary parameter space and the nature of the relativistic species present in the early universe. A particularly exciting frontier lies in secondary CMB anisotropies --- arising from gravitational lensing and scattering by free electrons in the late Universe --- which have so far been exploited primarily as astrophysical probes but hold significant untapped potential as windows into new physics, from non-gravitational dark matter interactions to non-Gaussian primordial signals. 

On the neutrino side, the combination of SO and galaxy surveys, such as DESI and Euclid, will bring the sum of neutrino masses within reach of a definitive cosmological measurement. This prospect is made all the more compelling by the current mild tension between cosmological data and the lower bound established by oscillation experiments~\cite{Craig:2016lyx,Green:2024xbb,Loverde:2024nfi,Lynch:2025ine}, which may yet point to new physics or overlooked systematics. 

Looking further ahead, once data from HERA and especially SKA become available, the 21-cm signal from cosmic dawn will open an entirely new observational window into the early Universe at redshifts currently beyond reach, with wide opportunities on the theoretical front to develop robust frameworks to interpret these observations in terms of fundamental physics.

These are just some of the directions in which the methods and results developed in this dissertation find their natural continuation --- as the interplay of theoretical insight and observational precision brings us ever closer to uncovering the missing physics that will complete our picture of the Universe.